% ============================================================
%  Chapter 1 -- Fields, 1.12-1.18
% ============================================================
\rsection{Fields}

\begin{signpost}[Why this section is here]
This is the second half of the vocabulary, and the less interesting half. An
ordered set gave Rudin the language for \emph{gaps}; a field gives him the
language for \emph{arithmetic}. $\R$ will be defined in 1.19 as an object having
both, plus one extra property, so both notions must be on the table first.

\smallskip
Rudin tells you outright how to read 1.14--1.18: they are ``some of the familiar
properties of $\Q$'' rederived from the axioms so that ``once we do this, we
will not need to do it again for the real numbers and for the complex numbers.''
That is the entire point. Nothing here will surprise you; the value is that
every one of these facts is now available for $\R$, for $\C$, and for any other
field, purchased once. Skim the proofs, note Definition 1.17, and move on ---
but see the warning after 1.18, because the ordering axioms are where students
later import assumptions that are not there.
\end{signpost}

\ritem{1.12}{Definition}
A \emph{field} is a set $F$ with two operations, called \emph{addition} and
\emph{multiplication}, which satisfy the following so-called ``field axioms''
(A), (M), and (D):

\smallskip
\noindent\textbf{(A) Axioms for addition}
\begin{enumerate}[label=(A\arabic*), leftmargin=3.1em, itemsep=2pt, topsep=3pt]
  \item If $x \in F$ and $y \in F$, then their sum $x+y$ is in $F$.
  \item Addition is commutative: $x + y = y + x$ for all $x, y \in F$.
  \item Addition is associative: $(x+y)+z = x+(y+z)$ for all $x,y,z \in F$.
  \item $F$ contains an element $0$ such that $0 + x = x$ for every $x \in F$.
  \item To every $x \in F$ corresponds an element $-x \in F$ such that
        \[ x + (-x) = 0. \]
\end{enumerate}

\smallskip
\noindent\textbf{(M) Axioms for multiplication}
\begin{enumerate}[label=(M\arabic*), leftmargin=3.1em, itemsep=2pt, topsep=3pt]
  \item If $x \in F$ and $y \in F$, then their product $xy$ is in
        $F$.\kw{\textbf{Source correction.} The retyped edition this book
        follows drops ``and $y \in F$'' here, and drops the opening ``If'' from
        (A1); its own foreword warns that new typos were likely introduced.
        Both clauses are restored above from Rudin's third edition. Without
        $y \in F$ the axiom does not say what it must --- that $F$ is
        \emph{closed} under multiplication.}
  \item Multiplication is commutative: $xy = yx$ for all $x, y \in F$.
  \item Multiplication is associative: $(xy)z = x(yz)$ for all $x,y,z \in F$.
  \item $F$ contains an element $1 \ne 0$ such that $1x = x$ for every
        $x \in F$.\kn{The clause $1 \ne 0$ is not pedantry. Without it the
        one-element set $\{0\}$ satisfies every other axiom, and then
        Proposition 1.16(a) would make every element equal to $0$. Excluding it
        is what keeps ``field'' from being vacuous.}
  \item If $x \in F$ and $x \ne 0$ then there exists an element $1/x \in F$ such
        that
        \[ x \cdot (1/x) = 1. \]
\end{enumerate}

\smallskip
\noindent\textbf{(D) The distributive law}
\[ x(y+z) = xy + xz \]
holds for all $x,y,z \in F$.\kd{(D) is the only axiom linking the two
operations, and every proof below that mixes $+$ and $\cdot$ must route through
it. Watch 1.16(a): the fact that $0x = 0$ --- which looks like arithmetic so
basic it needs no proof --- is derivable \emph{only} because (D) exists.}

\ritem{1.13}{Remarks}
\begin{enumerate}[label=(\alph*), leftmargin=2.2em, itemsep=5pt, topsep=4pt]
  \item One usually writes (in any field)
        \begin{gather*}
          x-y, \quad \frac{x}{y}, \quad x+y+z, \quad xyz,\\
          x^2, \quad x^3, \quad 2x, \quad 3x, \dots
        \end{gather*}
        in place of
        \begin{gather*}
          x+(-y), \quad x\cdot\left(\frac{1}{y}\right), \quad (x+y)+z,
          \quad (xy)z,\\
          xx, \quad xxx, \quad x+x, \quad x+x+x, \dots
        \end{gather*}
  \item The field axioms clearly hold in $\Q$, the set of all rational numbers,
        if addition and multiplication have their customary meaning. Thus $\Q$
        is a field.
  \item Although it is not our purpose to study fields (or any other algebraic
        structures) in detail, it is worthwhile to prove that some of the
        familiar properties of $\Q$ are consequences of the field axioms; once
        we do this, we will not need to do it again for the real numbers and for
        the complex numbers.
\end{enumerate}

\ritem{1.14}{Proposition}
\emph{The axioms for addition imply the following statements.}
\begin{enumerate}[label=(\alph*), leftmargin=2.2em, itemsep=1pt, topsep=3pt]
  \item \emph{If $x+y = x+z$ then $y = z$.}
  \item \emph{If $x+y = x$ then $y = 0$.}
  \item \emph{If $x+y = 0$ then $y = -x$.}
  \item \emph{$-(-x) = x$.}
\end{enumerate}
Statement (a) is a cancellation law. Note that (b) asserts the uniqueness of the
element whose existence is assumed in (A4), and that (c) does the same for
(A5).

\begin{proof}
If $x+y = x+z$, the axioms (A) give
\begin{align*}
  y &= 0 + y = (-x+x)+y = -x + (x+y)\\
    &= -x + (x+z) = (-x+x)+z = 0+z = z.
\end{align*}
This proves (a). Take $z = 0$ in (a) to obtain (b). Take $z = -x$ in (a) to
obtain (c). Since $-x+x = 0$, (c) (with $-x$ in place of $x$) gives (d).
\end{proof}

\begin{unpack}[The pattern of all four proofs]
Every step in that display is a single named axiom, and it is worth naming them
once: $0+y = y$ is (A4); $0 = -x+x$ is (A5) with (A2); the regrouping is (A3).
Nothing else is used --- in particular, no multiplication and no order.

The economy of (b), (c), (d) is the point. Once cancellation (a) is available,
the rest are not new arguments but \emph{substitutions} into it. Rudin's
``take $z = 0$'' and ``take $z = -x$'' are complete proofs, and the habit of
proving one general statement and specialising it repeatedly is one he will use
constantly.

\smallskip
How is the chain for (a) \emph{found}, rather than checked? Work from the
goal. You must delete $x$ from $x+y$; the only axiom that annihilates anything
is (A5), so $-x$ must enter; the only license to bring it next to $x$ is
regrouping, (A3); and once $-x+x$ becomes $0$, (A4) removes it. Each step is
forced --- the axioms offer exactly one tool per obstacle, so the proof writes
itself left to right once you know what must disappear. This is pattern 1 of
the strategy index: with abstract objects, unwind the definitions, because
there is nothing else to use.
\end{unpack}

\ritem{1.15}{Proposition}
\emph{The axioms for multiplication imply the following statements.}
\begin{enumerate}[label=(\alph*), leftmargin=2.2em, itemsep=1pt, topsep=3pt]
  \item \emph{If $x \ne 0$ and $xy = xz$ then $y = z$.}
  \item \emph{If $x \ne 0$ and $xy = x$ then $y = 1$.}
  \item \emph{If $x \ne 0$ and $xy = 1$ then $y = 1/x$.}
  \item \emph{If $x \ne 0$ then $1/(1/x) = x$.}
\end{enumerate}
The proof is so similar to that of Proposition 1.14 that we omit it.\kn{This is
Exercise 1.3. Copy the proof of 1.14 with $+ \to \cdot$, $0 \to 1$,
$-x \to 1/x$, and carry the hypothesis $x \ne 0$ along; the axioms (M1)--(M5)
match (A1)--(A5) exactly except for that restriction. The restriction is
unavoidable: (M5) grants inverses only to nonzero elements, which is precisely
why the multiplicative statements need a hypothesis the additive ones did not.}

\ritem{1.16}{Proposition}
\emph{The field axioms imply the following statements, for any $x,y,z \in F$.}
\begin{enumerate}[label=(\alph*), leftmargin=2.2em, itemsep=1pt, topsep=3pt]
  \item \emph{$0x = 0$.}
  \item \emph{If $x \ne 0$ and $y \ne 0$ then $xy \ne 0$.}
  \item \emph{$(-x)y = -(xy) = x(-y)$.}
  \item \emph{$(-x)(-y) = xy$.}
\end{enumerate}

\begin{proof}
$0x + 0x = (0+0)x = 0x$. Hence Proposition 1.14(b) implies that $0x = 0$, and
(a) holds.\kn{The trick in one line: manufacture an equation of the form
$u + u = u$, then quote 1.14(b) to conclude $u = 0$. How is it found? Ask what
the axioms \emph{know} about $0$: only (A4), $0 = 0+0$. So substitute that
and distribute --- the sole fact about $0$ meets the sole axiom linking $+$
to $\cdot$, and the proof is the only sentence one can write. First place (D)
does any work.} Next, assume $x \ne 0$, $y \ne 0$, but
$xy = 0$. Then
\[
  1 = \left(\frac1y\right)\left(\frac1x\right)xy
    = \left(\frac1y\right)\left(\frac1x\right)0 = 0,
\]
a contradiction. Thus (b) holds.

The first equality in (c) comes from
\[ (-x)y + xy = (-x+x)y = 0y = 0, \]
combined with Proposition 1.14(c); the other half of (c) is proved in the same
way. Finally,
\[ (-x)(-y) = -[x(-y)] = -[-(xy)] = xy \]
by (c) and Proposition 1.14(d).
\end{proof}

\begin{deeper}[What (b) actually says]
Part (b) is the statement that a field has \emph{no zero divisors}, and unlike
the rest of this proposition it is not a formality. There are perfectly good
rings where it fails: in the integers mod $6$, $2 \cdot 3 = 0$ with neither
factor zero. What rules that out here is (M5) --- the existence of inverses ---
and the proof shows exactly how: multiply by $1/x$ and $1/y$ and the supposed
$xy = 0$ collapses $1$ to $0$, contradicting $1 \ne 0$ in (M4).

This is why $\Z/6\Z$ is not a field while $\Z/5\Z$ is, and it is quietly
required every time a later argument divides by a nonzero quantity.
\end{deeper}

\ritem{1.17}{Definition}
An \emph{ordered field} is a field $F$ which is also an ordered set, such that
\begin{enumerate}[label=(\roman*), leftmargin=2.2em, itemsep=2pt, topsep=3pt]
  \item $x+y < x+z$ if $x,y,z \in F$ and $y < z$,
  \item $xy > 0$ if $x \in F$, $y \in F$, $x > 0$, and $y > 0$.
\end{enumerate}
If $x > 0$, we call $x$ \emph{positive}; if $x < 0$, $x$ is
\emph{negative}.\kn{Only two axioms, and they are the minimum needed to tie the
order to the arithmetic: translation preserves order, and positives are closed
under multiplication. Everything else you believe about inequalities ---
1.18(a)--(e), and the whole of Chapter 1's estimating --- is derived from these
two. Rudin is not listing convenient facts; he is exhibiting a complete basis.}

For example, $\Q$ is an ordered field.

All the familiar rules for working with inequalities apply in every ordered
field: Multiplication by positive [negative] quantities preserves [reverses]
inequalities, no square is negative, etc. The following proposition lists some of
these.

\ritem{1.18}{Proposition}
\emph{The following statements are true in every ordered field.}
\begin{enumerate}[label=(\alph*), leftmargin=2.2em, itemsep=1pt, topsep=3pt]
  \item \emph{If $x > 0$ then $-x < 0$, and vice versa.}
  \item \emph{If $x > 0$ and $y < z$ then $xy < xz$.}
  \item \emph{If $x < 0$ and $y < z$ then $xy > xz$.}
  \item \emph{If $x \ne 0$ then $x^2 > 0$. In particular, $1 > 0$.}
  \item \emph{If $0 < x < y$ then $0 < 1/y < 1/x$.}
\end{enumerate}

\begin{proof}
\begin{enumerate}[label=(\alph*), leftmargin=2.2em, itemsep=4pt, topsep=3pt]
  \item If $x > 0$ then $0 = -x+x > -x+0$, so that $-x < 0$. If $x < 0$ then
        $0 = -x+x < -x+0$, so that $-x > 0$. This proves (a).
  \item Since $z > y$, we have $z - y > y - y = 0$, hence $x(z-y) > 0$, and
        therefore
        \[ xz = x(z-y)+xy > 0 + xy = xy. \]
  \item By (a), (b), and Proposition 1.16(c),
        \[ -[x(z-y)] = (-x)(z-y) > 0, \]
        so that $x(z-y) < 0$, hence $xz < xy$.
  \item If $x > 0$, part (ii) of Definition 1.17 gives $x^2 > 0$. If $x < 0$,
        then $-x > 0$, hence $(-x)^2 > 0$. But $x^2 = (-x)^2$ by Proposition
        1.16(d). Since $1 = 1^2$, $1 > 0$.
  \item If $y > 0$ and $v \le 0$, then $yv \le 0$. But $y \cdot (1/y) = 1 > 0$.
        Hence $1/y > 0$. Likewise, $1/x > 0$. If we multiply both sides of the
        inequality $x < y$ by the positive quantity $(1/x)(1/y)$, we obtain
        $1/y < 1/x$.
\end{enumerate}
\end{proof}

\begin{unpack}[The one move used in every part]
Each proof converts an inequality into a statement about something being
positive, applies an axiom, and converts back. Concretely: $y < z$ becomes
$z - y > 0$ by (i) of 1.17, then (ii) applies, then (i) restores the inequality.

That is the whole technique. It is worth watching in (b), where it is cleanest:
$z>y \Rightarrow z-y>0 \Rightarrow x(z-y)>0 \Rightarrow xz > xy$. Part (c) is
(b) applied to $-x$, and (e) is (b) applied with the positive multiplier
$(1/x)(1/y)$. Only (d) does something different, splitting into cases on the
sign of $x$.
\end{unpack}

\begin{pitfall}[Part (d) is the source of most of Chapter 1's leverage]
$x^2 > 0$ for $x \ne 0$ looks too obvious to be worth stating. It is the reason
$\C$ cannot be ordered (Exercise 1.8: $-1 = i^2$ would have to be both positive
and negative), the reason $a^2 + b^2 > 0$ in the proof of (M5) for $\C$ at 1.25,
and the reason $|z| = (z\bar z)^{1/2}$ makes sense at 1.32.

The corollary $1 > 0$ deserves separate notice. It is not an axiom, and in a
field with no order there is no sense in which $1$ exceeds $0$. Here it is
\emph{derived}, from $1 = 1^2$ --- and once you have it, (i) of 1.17 gives
$0 < 1 < 2 < 3 < \cdots$, so every ordered field contains a copy of $\N$ in the
right order. That fact is used without comment in the proof of 1.20.
\end{pitfall}
